\chapter{Przecięcia, dualność i klasy charakterystyczne}
\label{ch:przeciecia-dualnosc-klasy}
\index{dualnosc Poincarego@Dualność Poincarégo}

Otoczenie tubularne z Rozdziału~\ref{ch:wlozenia-otoczenia}
pozwala odróżnić ruch \emph{wzdłuż} podrozmaitości
od ruchu \emph{poprzecznego}. Teraz przypiszemy
przecięciom liczby, a wiązkom -- klasy kohomologii.
W ten sposób geometryczne dane zanurzenia stają się
niezmiennikami, które można obliczać i porównywać.
To przygotowanie do teorii Morse'a, Smale'a i chirurgii;
w końcowej sekcji wykonamy też prostą operację
wycinania i sklejania rozmaitości.

Przez \emph{zamkniętą rozmaitość} rozumiemy w tym
rozdziale rozmaitość zwartą bez brzegu. W dowodzie
komórkowym będziemy używać triangulacji gładkiej
rozmaitości. Jej istnienie jest osobnym, głębszym
twierdzeniem o triangulacji; opisujemy dokładnie
argument dualności \emph{po} wybraniu takiej
triangulacji. Rozróżnienie tych dwóch kroków
zapobiega ukryciu twierdzenia o triangulacji
w rysunku.

\section{Kołańcuchy symplicjalne i iloczyn kubkowy}
\label{sec:kubek-16}

Kompleksy łańcuchowe i ich duale poznaliśmy
w Rozdziałach~\ref{ch:algebra-homologiczna}
i~\ref{ch:homologia-przestrzeni}. Dla
skończonej triangulacji $K$ i ciała $\mathbb F$
zapisujemy precyzyjnie
\[
 C^k(K;\mathbb F)
   =\operatorname{Hom}_{\mathbb F}
       (C_k(K;\mathbb F),\mathbb F),
 \qquad
 (\dd\varphi)(c)=\varphi(\partial c).
\]
\emph{Kołańcuch} $\varphi\in C^k$ jest więc
funkcją liniową \emph{na łańcuchach} stopnia $k$.
Nie jest kolejnym rodzajem sympleksu:
przypisuje każdemu łańcuchowi liczbę z
$\mathbb F$. To polska nazwa angielskiego
\emph{cochain}; łańcuch należy do $C_k$,
a kołańcuch do przestrzeni dualnej $C^k$.
Znak $\partial$ pozostaje zarezerwowany
dla łańcuchów, a $\dd$ dla kołańcuchów.
Równość $\dd^2=0$ wynika od razu z
$\partial^2=0$. Klasa w $H^k$ to
funkcja na cyklach, która nie zmienia
się po dodaniu brzegu. Nad ciałem
$H^k(K;\mathbb F)\cong
\operatorname{Hom}_{\mathbb F}(H_k(K;\mathbb F),\mathbb F)$:
można wybrać uzupełnienia baz jąder
i obrazów i zdualizować powstały
rozkład. Nad $\Z$ analogiczne zdanie
wymaga poprawki $\operatorname{Ext}$,
co było omawiane przy zmianie
współczynników w Rozdziale~\ref{ch:algebra-homologiczna}.
Pracę nad ciałem rozpoczniemy tu dlatego,
że łatwiej zobaczyć parowanie;
dualność Poincarégo ma również
wersję całkowitą.

\begin{definicja}[Iloczyn kubkowy]
\index{iloczyn kubkowy@Iloczyn kubkowy}
\label{def:kubek-16}
Ustalmy porządek wierzchołków $K$ i przemienny
pierścień współczynników $R$ (w szczególności $\Z$ lub ciało).
Dla $\alpha\in C^p(K;R)$ i $\beta\in C^q(K;R)$ definiujemy
na sympleksach o wierzchołkach zapisanych rosnąco
\[
(\alpha\smile\beta)
 [v_0,\ldots,v_{p+q}]
=\alpha[v_0,\ldots,v_p]\,
 \beta[v_p,\ldots,v_{p+q}].
\]
Na przeciwnie zorientowanym sympleksie przedłużamy ten
wzór liniowo ze znakiem. Obie ściany dzielą wierzchołek $v_p$.
Po przejściu do kohomologii otrzymamy
iloczyn $H^p\times H^q\to H^{p+q}$.
\end{definicja}

\begin{stwierdzenie}[Reguła Leibniza]
\label{prop:kubek-leibniz-16}
Dla kołańcuchów zachodzi
\[
 \dd(\alpha\smile\beta)
 =(\dd\alpha)\smile\beta
       +(-1)^p\alpha\smile(\dd\beta).
\]
Zatem iloczyn klas kohomologii jest
dobrze określony; jest łączny.
\end{stwierdzenie}
\begin{proof}
Obliczamy prawą stronę na $[v_0,\ldots,v_{p+q+1}]$.
W $(\dd\alpha)\smile\beta$ wyrazy z usunięciem
$v_i$ dla $0\leq i\leq p$ odpowiadają dokładnie
tym samym wyrazom w $\dd(\alpha\smile\beta)$.
W $(-1)^p\alpha\smile\dd\beta$ odpowiadają im
pozostałe wyrazy, o indeksach $p+1\leq i\leq p+q+1$.
Na prawej stronie są ponadto dwa dodatkowe wyrazy:
\[
 \bigl((-1)^{p+1}+(-1)^p\bigr)
 \alpha[v_0,\ldots,v_p]\,
 \beta[v_{p+1},\ldots,v_{p+q+1}]=0.
\]
Pochodzą one odpowiednio z ostatniej ściany
$\dd\alpha$ i pierwszej ściany $\dd\beta$.
To one się znoszą; żaden wyraz lewego kobrzegu
nie zostaje usunięty. Otrzymujemy regułę Leibniza.
Jeśli $\dd\alpha=\dd\beta=0$,
iloczyn jest kocyklem.
Jeśli $\alpha$ zastąpimy przez
$\alpha+\dd\gamma$, to zmiana
iloczynu wynosi
$\dd(\gamma\smile\beta)$
przy zamkniętej $\beta$;
przy zmianie $\beta$ na $\beta+\dd\eta$ różnica
wynosi $(-1)^p\dd(\alpha\smile\eta)$ przy zamkniętej $\alpha$.
Łączność widać bezpośrednio:
obie iteracje na sympleksie
stopnia $p+q+r$ dają iloczyn
wartości na trzech kolejnych
ścianach
$[v_0,\ldots,v_p]$,
$[v_p,\ldots,v_{p+q}]$ i
$[v_{p+q},\ldots,v_{p+q+r}]$.
\end{proof}

\begin{stwierdzenie}[Przemienność z właściwym znakiem]
\label{prop:kubek-graded-16}
W kohomologii
$\alpha\smile\beta
=(-1)^{pq}\beta\smile\alpha$
dla $\alpha\in H^p$,
$\beta\in H^q$.
\end{stwierdzenie}
\begin{proof}
Pracujemy najpierw nad $\Z$ i oznaczamy $C=C_*(K;\Z)$.
Na kompleksie tensorowym przyjmujemy
$\partial(a\otimes b)=\partial a\otimes b+(-1)^{|a|}a\otimes\partial b$.
Odwzorowanie Alexandera--Whitneya ma postać
\[
 A[v_0,\ldots,v_m]
 =\sum_{j=0}^m[v_0,\ldots,v_j]\otimes[v_j,\ldots,v_m].
\]
Jest odwzorowaniem łańcuchowym: po rozpisaniu brzegów
wyrazy wewnętrzne przy miejscu podziału znoszą się
tak jak w poprzednim dowodzie. Również
$T(a\otimes b)=(-1)^{|a||b|}b\otimes a$
jest odwzorowaniem łańcuchowym.

Skonstruujemy homotopię $H$ między $A$ i $TA$.
Dla wierzchołków oba odwzorowania są równe, więc
kładziemy $H=0$ w stopniu zero. Jeżeli $H$ jest
już określone na ścianach sympleksu $\sigma^m$,
łańcuch
\[
 z=TA(\sigma)-A(\sigma)-H(\partial\sigma)
\]
jest cyklem: stosujemy $\partial$, korzystamy
z łańcuchowości obu map i założenia indukcyjnego.
Cały $z$ leży w $C_*(\sigma)\otimes C_*(\sigma)$.
Ten kompleks ma zerową homologię w dodatnich stopniach:
stożkowanie sympleksu do jego wierzchołka jest
ściągnięciem łańcuchowym do $\Z$ w stopniu zero,
a tensorowanie tych ściągnięć daje to samo dla iloczynu.
Możemy więc wybrać $H(\sigma)$ z $\partial H(\sigma)=z$.
Indukcja daje $\partial H+H\partial=TA-A$.

Dla kocykli $\alpha,\beta$ kołańcuch
$\alpha\otimes\beta$ znika na brzegach kompleksu
tensorowego. Po ewaluacji ostatniej tożsamości
otrzymujemy, że
$(-1)^{pq}\beta\smile\alpha-\alpha\smile\beta$
jest kobrzegiem kołańcucha $(\alpha\otimes\beta)H$.
Dowód działa po zmianie współczynników na $R$.
Ta sama indukcja w ściągalnych nośnikach pokazuje
niezależność od porządku wierzchołków i wyboru
przybliżenia przekątnej. W połączeniu z aproksymacją
symplicjalną daje naturalność iloczynu w kohomologii.
\end{proof}

W przeciwieństwie do samego wymiaru
grup $H^k$, iloczyn kubkowy pamięta,
\emph{jak} różne cykle są względem siebie
położone. Ta dodatkowa informacja
wyjaśni później liczby przecięcia.
Oznaczenie $\smile$ odróżnia
iloczyn kołańcuchów od iloczynu
zewnętrznego form $\wedge$ z
Rozdziału~\ref{ch:formy-de-rham}.

\begin{definicja}[Iloczyn kapowy i konwencja znaków]
\index{iloczyn kapowy i konwencja znakow@Iloczyn kapowy i konwencja znaków}
\label{def:kapowy-16}
Dla $m$-sympleksu $c=[v_0,\ldots,v_m]$ i kołańcucha
$\alpha$ stopnia $k\leq m$ kładziemy
\[
 c\frown\alpha=\alpha[v_0,\ldots,v_k]\,[v_k,\ldots,v_m].
\]
Wzór przedłużamy liniowo. Możemy wziąć całkowity łańcuch
i kołańcuch o wartościach w dowolnej grupie abelowej $G$;
wtedy wynik jest łańcuchem z $G$. Bezpośrednie rozpisanie
ścian (z tym samym skróceniem dwóch wyrazów przy podziale)
daje
\begin{equation}\label{eq:kapowy-brzeg-16}
 \partial(c\frown\alpha)=(-1)^k
  \bigl((\partial c)\frown\alpha-c\frown\dd\alpha\bigr).
\end{equation}
Dlatego cykl i kocykl wyznaczają klasę homologii,
niezależną od reprezentantów. Dla współczynników w $R$
mamy ponadto, prosto ze wzorów na ściany,
\[
 \langle\beta,c\frown\alpha\rangle
 =\langle\alpha\smile\beta,c\rangle.
\]
Ta kolejność czynników będzie obowiązywała w całym rozdziale.
\end{definicja}

\section{Orientacja jako klasa homologii}
\label{sec:klasa-fundamentalna-16}

W Stwierdzeniu~\ref{prop:lokalna-homologia}
ustaliliśmy, że
$H_n(M,M\setminus\{x\};\Z)\cong\Z$.
Orientacja oznacza spójny wybór
generatorów tych lokalnych grup.
Teraz skleimy je w jedną klasę
globalną.

\begin{twierdzenie}[Klasa fundamentalna]
\label{thm:klasa-fundamentalna-16}
Niech $M^n$ będzie spójną, zamkniętą,
zorientowaną rozmaitością. Istnieje
jedyna klasa
$[M]\in H_n(M;\Z)$, która w
$H_n(M,M\setminus\{x\};\Z)$
przechodzi w wybrany lokalny
generator dla każdego $x$.
W szczególności $H_n(M;\Z)\cong\Z$.
Dla zwartej zorientowanej rozmaitości z brzegiem
otrzymujemy zamiast tego
$[M,\partial M]\in H_n(M,\partial M;\Z)$.
Bez orientacji zawsze istnieje
klasa fundamentalna ze
współczynnikami $\mathbb F_2=\Z/2$.
\end{twierdzenie}
\begin{proof}
Wybierzmy triangulację $K$ zgodną
ze strukturą gładką. Orientacja $M$
wyznacza orientację każdego
$n$-sympleksu. Niech
$c=\sum_{\sigma^n\in K}\sigma^n$
oznacza sumę tak zorientowanych
sympleksów (w zapisie łańcuchowym
znaki są już uwzględnione w ich
orientacjach). Każda wewnętrzna
$(n-1)$-ściana należy dokładnie
do dwóch $n$-sympleksów.
Orientacje indukowane na ich
wspólnej ścianie są przeciwne,
więc odpowiadające współczynniki
w $\partial c$ się znoszą.
W przypadku $\partial M=\varnothing$
otrzymujemy $\partial c=0$.
W przypadku brzegu pozostają tylko
ściany leżące w $\partial M$,
więc $\partial c\in C_{n-1}(\partial M)$:
$c$ jest cyklem względnym.

W punkcie $x$ we wnętrzu
$n$-sympleksu obraz $c$ w
grupie lokalnej jest generatorem
o znaku określonym orientacją.
Jeśli $x$ leży na ścianie,
mała gwiazda wokół $x$ jest
kulą złożoną z przyległych
sympleksów; znoszenie się
wewnętrznych ścian pokazuje,
że ta sama suma daje generator
również tam. W przypadku brzegu powyższy warunek lokalnego
generatora nakładamy na punkty \emph{wnętrza}.
Dla $x\in\partial M$ zwykła grupa
$H_n(M,M\setminus\{x\};\Z)$ jest zerowa:
półkula i półkula z usuniętym środkiem płaskiej ściany
są ściągalne (dla $n\geq1$). Nie daje więc lokalnej
orientacji brzegu. Zamiast tego brzeg klasy względnej
spełnia $\partial[M,\partial M]=[\partial M]$
przy orientacji „zewnętrzna normalna najpierw”.

Wyjaśnijmy jedyność w przypadku
zamkniętym. W każdym $n$-cyklu
$z=\sum a_\sigma\sigma$ warunek
$\partial z=0$ przy wspólnej
$(n-1)$-ścianie mówi, że
współczynniki $a_\sigma$ po obu
jej stronach są równe
(po uwzględnieniu orientacji).
Graf sąsiedztwa najwyższych
sympleksów jest spójny,
gdy $M$ jest spójna i $n\geq1$:
droga między ich wnętrzami
może po małej perturbacji
omijać szkielet wymiaru
$n-2$ i kolejno przecinać
ściany kowymiaru jeden.
Wszystkie współczynniki są
więc jednakowe. Ponieważ
$C_{n+1}(K)=0$, nie ma
nietrywialnych symplicjalnych
brzegów w stopniu $n$.
Stąd $H_n(M;\Z)=\Z[c]$. Dla $n=0$ spójna
zamknięta rozmaitość jest punktem i wynik jest bezpośredni.
W przypadku względnym to samo porównanie współczynników
przez ściany wewnętrzne daje jedyność klasy względnej.
Homologia symplicjalna jest
zgodna z singularną na
triangulacji, jak pokazano
w Rozdziale~\ref{ch:homologia-przestrzeni}.
Jeśli porzucimy znaki i pracujemy
nad $\mathbb F_2$, dwie ściany
sumują się do $0$ bez wybierania
orientacji. Ten sam argument
daje klasę modulo $2$.
\end{proof}

\begin{przyklad}
Na $S^2$ suma zorientowanych
trójkątów daje $[S^2]$.
Na torusie analogiczna suma
daje $[T^2]$. Dla wstęgi
Möbiusa nie istnieje
klasa $[M,\partial M]$ nad $\Z$
o lokalnych współczynnikach
równych $+1$ w każdym punkcie:
po obejściu rdzenia orientacja
zmienia znak. Klasa modulo $2$
istnieje. Poprawną wersję
całkowitą uzyskamy z lokalnymi
współczynnikami w
Sekcji~\ref{sec:lokalne-wsp-16}.
\end{przyklad}

\section{Dualność Poincarégo przez komórki dualne}
\label{sec:dualnosc-komorkowa-16}

Niech $\sigma$ będzie $k$-sympleksem
triangulacji $K$ zamkniętej
$n$-rozmaitości. Po podziale
barycentrycznym bierzemy wszystkie
małe sympleksy odpowiadające
łańcuchom ścian
$\sigma=\sigma_k\subsetneq
\sigma_{k+1}\subsetneq\cdots
\subsetneq\sigma_n$.
Ich suma, z odpowiednimi
ścianami, tworzy
$(n-k)$-wymiarową
\emph{komórkę dualną} $D(\sigma)$.
Komórka $D(\sigma)$ przecina
$\sigma$ w barycentrum i nie
przecina wnętrza innych
$k$-sympleksów. W triangulacji zgodnej ze strukturą gładką link
$\sigma$ jest sferą PL (dla najwyższych sympleksów
komórki dualne są punktami), więc
$D(\sigma)$ jest domkniętym
dyskiem, a jego brzeg składa
się z komórek dualnych do
$(k+1)$-sympleksów
zawierających $\sigma$.

\begin{figure}[htbp]
\centering
\begin{tikzpicture}[>=Latex,font=\small,line cap=round,scale=1.08]
 \coordinate (o) at (0,0);
 \coordinate (a) at (-2,-1.65);
 \coordinate (b) at (2,-1.65);
 \coordinate (c) at (2,1.65);
 \coordinate (d) at (-2,1.65);
 \fill[manifoldgold!18]
  (-1,-.825)--(0,-1.10)--(1,-.825)--
  ({4/3},0)--(1,.825)--(0,1.10)--
  (-1,.825)--({-4/3},0)--cycle;
 \draw[manifoldblue!65,thick] (a)--(b)--(c)--(d)--cycle;
 \foreach \v in {a,b,c,d}
  \draw[manifoldblue!65,thick] (o)--(\v);
 \draw[manifoldgold!85!black,very thick]
  (-1,-.825)--(0,-1.10)--(1,-.825)--
  ({4/3},0)--(1,.825)--(0,1.10)--
  (-1,.825)--({-4/3},0)--cycle;
 \foreach \x/\y in {0/-1.10,0/1.10}
  \fill[manifoldgold!85!black] (\x,\y) circle (1.4pt);
 \fill[manifoldgold!85!black] ({4/3},0) circle (1.4pt);
 \fill[manifoldgold!85!black] ({-4/3},0) circle (1.4pt);
 \fill[manifoldred] (o) circle (2.5pt)
  node[below right=2pt] {$\sigma^0$};
 \node[manifoldgold!65!black] at (0,.65) {$D(\sigma^0)$};
 \node[manifoldblue,anchor=south] at (0,1.78) {triangulacja};
\end{tikzpicture}
\caption{Niebieskie trójkąty stanowią fragment
triangulacji powierzchni. Złoty dysk jest
komórką dualną do środkowego wierzchołka.
Analogicznie krawędź ma dualną krawędź
przecinającą ją poprzecznie.}
\label{fig:dualna-komorka-16}
\end{figure}

\begin{twierdzenie}[Dualność Poincarégo]
\label{thm:poincare-dualnosc-16}
Dla zamkniętej zorientowanej
rozmaitości $M^n$ i dowolnych
abelowych grup współczynników $G$ mnożenie
kapowe przez klasę fundamentalną
jest izomorfizmem
\[
 \operatorname{PD}:H^k(M;G)
       \stackrel{\cong}{\longrightarrow}
       H_{n-k}(M;G),\qquad
 \alpha\longmapsto[M]\frown\alpha.
\]
Dla zwartej zorientowanej
rozmaitości z brzegiem
odpowiednie izomorfizmy to
\[
 H^k(M;G)\cong H_{n-k}(M,\partial M;G),
 \qquad
 H^k(M,\partial M;G)\cong H_{n-k}(M;G).
\]
Przy $G=\mathbb F_2$ orientacja
nie jest potrzebna.
\end{twierdzenie}
\begin{proof}
Najpierw niech $\partial M=\varnothing$.
Wybieramy dualny rozkład komórkowy
$K^\vee$. Ustalmy orientacje sympleksów. Dla $G=\Z$ niech
$\sigma^\#$ będzie kołańcuchem równym $1$ na $\sigma$
i $0$ na pozostałych sympleksach tego wymiaru.
Orientujemy $D(\sigma)$ tak, aby kolejność
$(T\sigma,TD(\sigma))$ dawała orientację $M$.
Dla dowolnej grupy abelowej przenosimy po prostu
współczynnik $g\in G$ przy $\sigma^\#$ na tę komórkę.
Daje to izomorfizmy grup
\[
 \mathcal D_k:C^k(K;G)\longrightarrow C_{n-k}(K^\vee;G).
\]
Rachunek brzegu ma konkretny znak. Jeśli $\tau^{k+1}$
zawiera $\sigma^k$ i $u$ wskazuje od $\sigma$ do wnętrza
$\tau$, to konwencja brzegu daje
$\operatorname{or}(\tau)=[\tau:\sigma](-u,\operatorname{or}(\sigma))$.
Po przeniesieniu $u$ za $k$ wektorów otrzymujemy
$[\tau:\sigma](-1)^{k+1}(\operatorname{or}(\sigma),u)$.
W komórce dualnej kierunek $u$ wskazuje na zewnątrz
ku ścianie $D(\tau)$. Porównanie orientacji $M$ daje
\begin{equation}\label{eq:dualny-lancuch-16}
 \partial\mathcal D_k=(-1)^{k+1}\mathcal D_{k+1}\dd.
\end{equation}
Mimo znaku izomorfizmy $\mathcal D_k$ przenoszą
kocykle na cykle i kobrzegi na brzegi, zatem dają
$H^k(K;G)\cong H_{n-k}(K^\vee;G)$.
Jeśli chcemy dosłownego izomorfizmu kompleksów
z różniczkami $\dd$ i $\partial$, należy użyć
$(-1)^{k(k+1)/2}\mathcal D_k$. Zachowujemy jednak
nieprzeskalowane $\mathcal D_k$, zgodne ze znakiem
iloczynu kapowego w~\eqref{eq:kapowy-brzeg-16}.

Wyjaśnijmy utożsamienie z $[M]\frown\alpha$.
W podziale barycentrycznym obie konstrukcje dla
$\sigma^\#$ mają nośnik w zamkniętej gwieździe $\sigma$.
W stopniu $k=n$ dają punkty z tym samym współczynnikiem
orientacji. Różnicę można wypełnić drogą w tej gwieździe.
Schodząc kolejno z $k$ do $k-1$, z obu wzorów na brzeg
wynika, że różnica po odjęciu już zbudowanej homotopii
na kościanach jest cyklem w tej samej gwieździe.
Gwiazda jest ściągalna, więc cykl można wypełnić.
To indukcja analogiczna do dowodu przemienności
iloczynu kubkowego; daje homotopię po sprowadzeniu
obu konstrukcji do wspólnego podziału. Wyznaczają
więc tę samą mapę na kohomologii, czyli $\operatorname{PD}$.
Oba rozkłady komórkowe obliczają homologię $M$.

Gdy $\partial M\ne\varnothing$,
dualna komórka przy brzegu
ma dodatkową ścianę
\emph{leżącą w brzegu}.
Po przejściu do łańcuchów
względnych ta ściana znika;
daje to pierwszy z podanych
izomorfizmów. Z kolei
kołańcuchy względne znikają
na sympleksach brzegowych,
więc ich dualne komórki
nie mają obciętej ściany:
tworzą podkompleks dualny $L$ z komórek
$D(\sigma)$ dla $\sigma\not\subset\partial M$.
Ten podkompleks oblicza homologię $M$.
Istotnie, $M\setminus\partial M$ cofa się na $L$:
w każdym barycentrycznym sympleksie zerujemy wagi
wierzchołków odpowiadających ścianom brzegowym
i normalizujemy pozostałe. Ich suma jest dodatnia
poza brzegiem, a homotopia przez liniową zmianę wag
jest zgodna na ścianach. Włączenie wnętrza do $M$
jest równoważnością homotopijną, bo kołnierz pozwala
odsunąć brzeg do wnętrza. Stąd drugi izomorfizm
z homologii bezwzględnej.
W każdym przypadku
sprawdzenie~\eqref{eq:dualny-lancuch-16}
jest tym samym lokalnym
rachunkiem incydencji;
kołnierz z
Twierdzenia~\ref{thm:kolnierz-15}
zapewnia standardowy model
półprzestrzenny przy brzegu.
Nad $\mathbb F_2$ wszystkie
znaki orientacji znikają.
\end{proof}

\begin{wniosek}[Parowanie i charakterystyka Eulera]
\label{cor:pd-parowanie-16}
Dla zamkniętej zorientowanej
$M^n$ i ciała $\mathbb F$
parowanie
\[
 H^k(M;\mathbb F)\times
 H^{n-k}(M;\mathbb F)
 \longrightarrow\mathbb F,\qquad
 (\alpha,\beta)\longmapsto
      \langle\alpha\smile\beta,[M]\rangle
\]
jest niezdegenerowane. Zatem
$b_k=b_{n-k}$. Jeśli $n$
jest nieparzyste, $\chi(M)=0$.
\end{wniosek}
\begin{proof}
Dualność identyfikuje
$H^k$ z $H_{n-k}$.
Nad ciałem każda niezerowa
klasa homologii jest wykrywana
przez pewną klasę kohomologii
z $H^{n-k}$, bo dual przestrzeni
wektorowej rozdziela punkty.
Definicja iloczynu kapowego
i kubkowego utożsamia tę
ewaluację z podaną liczbą.
Stąd parowanie jest
niezdegenerowane i w skończonym
wymiarze obie strony mają
równe wymiary. Jeśli $n$
jest nieparzyste, w sumie
$\chi(M)=\sum_k(-1)^kb_k$
wyrazy dla $k$ i $n-k$
mają przeciwne znaki,
więc sumują się do zera.
\end{proof}

\begin{przyklad}[Torus i iloczyn dwóch sfer]
Na $T^2$ klasy
$a,b\in H^1(T^2;\Z)$,
odpowiadające dwóm kierunkom,
spełniają
$\langle a\smile b,[T^2]\rangle=1$
po ustaleniu orientacji.
Na $S^2\times S^2$ dwie
klasy $A,B\in H_2(-;\Z)$
reprezentowane przez
$S^2\times\{*\}$ i
$\{*\}\times S^2$
mają macierz przecięć
$\begin{psmallmatrix}0&1\\1&0\end{psmallmatrix}$.
Samych rang $H_2\cong\Z^2$
nie wystarcza do odtworzenia
tej macierzy.
\end{przyklad}

\section{Liczba przecięcia: znak i niezmienniczość}
\label{sec:przeciecia-16}

Niech $A^p,B^{n-p}\subset M^n$ będą
zamkniętymi zorientowanymi
podrozmaitościami zamkniętej
zorientowanej $M$. Po małej
perturbacji, możliwej dzięki
Twierdzeniu~\ref{thm:gestosc-transwersalnych},
możemy założyć $A\pitchfork B$.
Przecięcie ma wymiar zero,
a ze zwartości jest skończone.

\begin{definicja}[Znak i liczba przecięcia]
\index{znak i liczba przeciecia@Znak i liczba przecięcia}
\label{def:liczba-przeciecia-16}
Dla $x\in A\cap B$ kładziemy
$\epsilon_x=+1$, jeżeli
uporządkowana baza
$T_xA$ \emph{a następnie}
$T_xB$ daje orientację $T_xM$,
i $\epsilon_x=-1$ w przeciwnym
przypadku. \emph{Liczba
przecięcia} to
\[
 I(A,B)=\sum_{x\in A\cap B}\epsilon_x.
\]
Bez orientacji liczymy
$\#(A\cap B)\pmod 2$.
\end{definicja}

\begin{figure}[htbp]
\centering
\begin{tikzpicture}[>=Latex,font=\small,line cap=round,scale=1.04]
 \fill[manifoldblue!5] (-2.1,-1.65) rectangle (2.1,1.65);
 \draw[manifoldblue!45,thick]
   (-2.1,-1.65) rectangle (2.1,1.65);
 \draw[manifoldblue,very thick]
   (-2.1,0)--(2.1,0);
 \draw[manifoldblue,very thick,->]
   (1.25,0)--(1.75,0);
 \draw[manifoldgold!85!black,very thick]
   (0,-1.65)--(0,1.65);
 \draw[manifoldgold!85!black,very thick,->]
   (0,.60)--(0,1.10);
 \fill[manifoldred] (0,0) circle (2.6pt);
 \node[manifoldred,anchor=south west] at (.10,.08) {$x,\;\epsilon_x=+1$};
 \node[manifoldblue,anchor=north] at (1.35,-.12) {$A$};
 \node[manifoldgold!75!black,anchor=west] at (.12,1.15) {$B$};
 \node[anchor=south] at (0,1.72)
       {przeciwległe brzegi są sklejone};
\end{tikzpicture}
\caption{Dwa podstawowe okręgi torusa
przedstawionego jako kwadrat z parami
sklejonych boków. W punkcie $x$ kierunek
poziomy, po nim pionowy, daje dodatnią
orientację. Liczba przecięcia wynosi $1$.}
\label{fig:przeciecie-torus-16}
\end{figure}

\begin{twierdzenie}[Przecięcie i iloczyn kubkowy]
\label{thm:przeciecie-kubek-16}
Niech $\alpha_A=\operatorname{PD}^{-1}
(i_{A*}[A])\in H^{n-p}(M;\Z)$
i $\alpha_B=\operatorname{PD}^{-1}
(i_{B*}[B])\in H^p(M;\Z)$.
Przy konwencji orientowania przecięcia
z Definicji~\ref{def:liczba-przeciecia-16}
zachodzi
\begin{equation}\label{eq:przeciecie-kubek-16}
 I(A,B)=
 \langle\alpha_A\smile\alpha_B,[M]\rangle.
\end{equation}
Liczba zależy więc tylko od klas
homologii obu podrozmaitości;
jest niezmienna przy izotopii.
Jeżeli jest niezerowa, żadna
perturbacja w tych klasach nie
może rozdzielić $A$ i $B$.
\end{twierdzenie}
\begin{proof}
Użyjemy klasy Thoma z Twierdzenia~\ref{thm:euler-thom-16}.
Jej zamieszczony niżej dowód jest komórkowy i nie
korzysta z obecnego twierdzenia ani z klas $w_i$.
Orientujemy normalne warunkiem „normalna, potem styczna”.
Wycięcie w otoczeniu tubularnym daje klasy
$u_A\in H^{n-p}(M,M\setminus A;\Z)$ i
$u_B\in H^p(M,M\setminus B;\Z)$.
Ich obrazy w kohomologii bezwzględnej są $\alpha_A,\alpha_B$:
lokalnie oba opisy przypisują $+1$ dodatnio zorientowanemu
normalnemu dyskowi, co jest konstrukcją dualnej klasy.
Względny iloczyn ma nośnik na przecięciu, tzn.
\[
 u_A\smile u_B\in H^n(M,M\setminus(A\cap B);\Z).
\]
Wynika to z wersji względnej wzoru kubkowego:
po dostatecznym podziale sympleksy poza przecięciem
leżą w jednym z dwóch otwartych dopełnień,
gdzie odpowiedni czynnik znika.
Wycięcie rozkłada ostatnią grupę na sumę grup
$H^n(B_x,B_x\setminus\{x\};\Z)\cong\Z$
dla małych rozłącznych kul wokół punktów przecięcia.
Ewaluacja na $[M]$ jest sumą tych lokalnych liczb.
Jest to argument całkowity, nie tylko rachunek form
rzeczywistych, który mógłby zgubić torsję.

Oto kontrola znaku.
Niech $q=n-p$, a
$(a_1,\ldots,a_p)$
i $(b_1,\ldots,b_q)$
będą dodatnimi
bazami $T_xA$ i $T_xB$.
Załóżmy najpierw,
że $(a,b)$ orientuje
$T_xM$ dodatnio.
Normalną do $A$
orientujemy warunkiem
$(\nu_A,a)>0$,
a normalną do $B$
warunkiem
$(\nu_B,b)>0$.
Wówczas odpowiadająca
pierwszej normalnej
$q$-forma $\omega_A$
ma na bloku $b$
wartość $(-1)^{pq}$,
a $p$-forma $\omega_B$
drugiej normalnej
ma na bloku $a$
wartość $+1$.
We wzorze na
$(\omega_A\wedge\omega_B)(a,b)$
przestawienie bloku
$b$ przed $a$
daje drugi czynnik
$(-1)^{pq}$.
Iloczyn dwóch znaków
wynosi $+1$.
Odwrócenie orientacji
jednego z $A,B$
zmienia zarówno
$\epsilon_x$, jak
i tę wartość na $-1$.
Zatem lokalna ewaluacja
iloczynu jest $\epsilon_x$.
Po zsumowaniu po rozłącznych
małych otoczeniach punktów
przecięcia otrzymujemy
\eqref{eq:przeciecie-kubek-16}.
W powyższym rachunku formy $\omega_A,\omega_B$
opisują orientacje liniowych bloków; całkowite
generatory są ustalone przez klasy Thoma.
Poza tymi kulami nie ma wkładu dzięki względnemu
nośnikowi iloczynu. Ten sam dowód modulo $2$
działa bez orientacji. Dla transwersalnych przecięć
dodatniego wymiaru identyfikuje iloczyn klas Thoma
z klasą Thoma sumy normalnych; daje analogiczną
interpretację iloczynu jako klasy dualnej do przecięcia.

Prawa strona zależy jedynie
od $i_{A*}[A]$ i $i_{B*}[B]$,
ponieważ dualność Poincarégo
i iloczyn kubkowy są
określone na klasach.
Izotopijne zanurzenia
indukują homotopijne
odwzorowania w $M$,
więc dają tę samą
klasę homologii.
Gdyby obrazy dało się
rozłączyć, suma po
pustym przecięciu
wynosiłaby $0$.
\end{proof}

\begin{przyklad}[Krzywe na torusie]
W standardowej bazie
$H_1(T^2;\Z)=\Z a\oplus\Z b$
krzywe o klasach
$pa+qb$ i $ra+sb$
mają przecięcie algebraiczne
\[
 I\bigl((p,q),(r,s)\bigr)
 =\det\begin{pmatrix}p&r\\q&s\end{pmatrix}
 =ps-qr.
\]
Równość wynika z dwuliniowości
iloczynu i z
$I(a,b)=1$, $I(b,a)=-1$,
$I(a,a)=I(b,b)=0$.
Na przykład $(1,1)$ i $(1,-1)$
mają liczbę przecięcia $-2$:
co najmniej dwa przecięcia,
licząc bez znaków, są
nieuniknione w tych klasach.
Liczba $0$ nie gwarantuje
rozłączności konkretnych
przedstawicieli: dodatnie
i ujemne przecięcia mogą
się znosić.
\end{przyklad}

\begin{uwaga}[Samoprzecięcie]
Dwie kopie tej samej
podrozmaitości nie są
transwersalne, gdy się
pokrywają. Jedną z nich
przesuwamy w małym
otoczeniu tubularnym
za pomocą przekroju
wiązki normalnej.
Jeżeli przekrój ma zera,
przesunięta kopia
przecina pierwotną
właśnie w tych zerach.
Ta obserwacja prowadzi
do klasy Eulera
w Sekcji~\ref{sec:klasy-eulera-16}.
\end{uwaga}

\section{Lokalne współczynniki i brak orientacji}
\label{sec:lokalne-wsp-16}

Zwykła grupa współczynników
$\Z$ przyjmuje, że liczba
$+1$ oznacza to samo w
każdym punkcie. Na rozmaitości
nieorientowalnej lokalny
generator orientacji wraca
po pewnej pętli jako $-1$.
Trzeba więc przenosić
współczynniki wraz z drogą.

\begin{definicja}[System lokalny i system orientacji]
\index{system lokalny i system orientacji@System lokalny i system orientacji}
\label{def:system-lokalny-16}
\emph{System lokalny} grup abelowych
$\mathcal L$ przypisuje punktowi
$x$ grupę $\mathcal L_x$,
a drodze $\gamma:x\to y$
izomorfizm
$T_\gamma:\mathcal L_x\to\mathcal L_y$,
niezmienny przy homotopii
drogi z ustalonymi końcami
i zgodny ze składaniem dróg.
Transport wzdłuż drogi stałej jest identycznością.
Dla rozmaitości bez brzegu \emph{system orientacji}
$\mathcal O_M$ ma włókno
\[
 (\mathcal O_M)_x
  =H_n(M,M\setminus\{x\};\Z)
  \cong\Z;
\]
transport lokalnej orientacji
wzdłuż drogi jest mnożeniem
przez $+1$ albo $-1$. Przy brzegu definiujemy
system przez orientacje $T_xM$, równoważnie przez
przedłużenie systemu z wnętrza w kołnierzu.
Nie stosujemy tam powyższego wzoru na lokalną homologię.
\end{definicja}

W łańcuchach z $\mathcal L$
sympleks singularny ma
współczynnik w grupie
nad jego pierwszym
wierzchołkiem. Przy operatorze
brzegu współczynnik ściany
pomijającej pierwszy
wierzchołek przenosimy
wzdłuż pierwszej krawędzi
do nowego wierzchołka;
dla pozostałych ścian
punkt bazowy się nie
zmienia. Znaki przemienne
są takie jak w zwykłym
operatorze $\partial$.
Dwie drogi transportu
przy dwukrotnym usunięciu
wierzchołków są homotopijne
w obrębie sympleksu,
więc odpowiadające
wyrazy w $\partial^2$
znoszą się parami.
Tak powstaje
$H_*(M;\mathcal L)$.

\begin{przyklad}[Skręcone współczynniki na okręgu]
Niech $\mathcal L$ nad $S^1$
ma włókno $\Z$ i transport
po pełnym okrążeniu
równy $-1$. W komórkowym
modelu okręgu jest
jedna $0$-komórka
i jedna $1$-komórka.
Przy zwykłych
współczynnikach końce
krawędzi dają
$\partial_1=1-1=0$.
Tutaj jeden koniec
wraca ze znakiem
$-1$, więc
$\partial_1=(-1)-1=-2$.
Stąd
\[
 H_1(S^1;\mathcal L)=0,\qquad
 H_0(S^1;\mathcal L)=\Z/2.
\]
Po przejściu do
$\mathbb F_2$ macierz
$-2$ staje się zerem
i dostajemy zwykłe
grupy okręgu.
\end{przyklad}

\begin{twierdzenie}[Dualność ze skręconymi współczynnikami]
\label{thm:pd-skrecenie-16}
Dla spójnej zamkniętej
$n$-rozmaitości, także
nieorientowalnej, istnieje
kanoniczna skręcona klasa
fundamentalna
$[M]_{\mathcal O}\in
H_n(M;\mathcal O_M)$.
Dla stałej abelowej grupy
współczynników $G$
otrzymujemy izomorfizm
\[
 H^k(M;G)\cong
 H_{n-k}(M;\mathcal O_M\otimes G).
\]
Dla zwartej rozmaitości z brzegiem używamy
klasy $[M,\partial M]_{\mathcal O}$ i otrzymujemy
obie wersje względne:
\[
 H^k(M;G)\cong H_{n-k}(M,\partial M;\mathcal O_M\otimes G),
 \qquad
 H^k(M,\partial M;G)\cong H_{n-k}(M;\mathcal O_M\otimes G).
\]
Po zredukowaniu
współczynników modulo $2$
system orientacji staje
się stały.
\end{twierdzenie}
\begin{proof}
W każdej małej mapie
wybieramy generator
lokalnej homologii.
Gdy przechodzimy
do sąsiedniej mapy,
ewentualna zmiana
orientacji o $-1$
zmienia jednocześnie
współczynnik w
$\mathcal O_M$ o $-1$.
Iloczyn tych dwóch
znaków jest $+1$.
Suma wszystkich
najwyższych sympleksów
z tak dobranymi
lokalnymi współczynnikami
ma zatem zerowy brzeg
na każdej wspólnej
ścianie i daje
$[M]_{\mathcal O}$.
W małej kuli jej obraz
jest lokalnym generatorem.

W dowodzie
Twierdzenia~\ref{thm:poincare-dualnosc-16}
orientacja $M$ była
potrzebna do orientowania
komórek $D(\sigma)$.
Teraz orientujemy
$D(\sigma)$ za pomocą
lokalnego generatora
$(\mathcal O_M)_x$.
Przejście do sąsiedniej
komórki wnosi dokładnie
znak transportu
w $\mathcal O_M$.
Rachunek incydencji
\eqref{eq:dualny-lancuch-16}
pozostaje więc prawdziwy
z tymi współczynnikami.
Znowu dostajemy
izomorfizmy przenoszące jądra i obrazy
(ze znakiem zależnym od stopnia), a więc ich homologie.
Na brzegu przechodzimy
do łańcuchów względnych
tak samo jak w dowodzie
wersji zorientowanej.
W $\mathbb F_2$
mamy $-1=+1$,
więc monodromia znika.
\end{proof}

\begin{przyklad}[Wstęga Möbiusa i płaszczyzna rzutowa]
Dla wstęgi Möbiusa $W$
nie istnieje zwykła
całkowita klasa
$[W,\partial W]$
generująca orientację
we wszystkich punktach,
ale istnieje
$[W,\partial W]_{\mathcal O}
\in H_2(W,\partial W;\mathcal O_W)
\cong\Z$.
Dla $\R P^2$ mamy
$H_2(\R P^2;\Z)=0$,
za to
$H_2(\R P^2;\mathcal O)\cong\Z$
i $H_2(\R P^2;\mathbb F_2)
\cong\mathbb F_2$.
Lokalne współczynniki
odzyskują orientację
bez twierdzenia,
że przestrzeń nagle
stała się orientowalna.
\end{przyklad}

\section{Wiązki nad sferami i funkcje sklejające}
\label{sec:clutching-16}

Sferę $S^m$ rozkładamy na dwa
domknięte dyski $D^m_+$ i $D^m_-$,
sklejone wzdłuż równika
$S^{m-1}$. Wiązka wektorowa
nad każdym dyskiem jest
trywialna, ponieważ dysk
jest ściągalny. Cała informacja
o wiązce leży więc w
\emph{funkcji sklejającej}
na równiku.

\begin{twierdzenie}[Klasyfikacja przez równik]
\label{thm:clutching-16}
Dla $m\geq2$ klasy izomorfizmu zorientowanych
rzeczywistych wiązek
rzędu $r\geq1$ nad $S^m$
odpowiadają klasom homotopii
\[
 [S^{m-1},SO(r)]
 \cong\pi_{m-1}(SO(r)).
\]
Funkcja $g:S^{m-1}\to SO(r)$
skleja dwie wiązki
trywialne przez
$(x,v)_+\sim(x,g(x)v)_-$.
Izomorfizmy w tej klasyfikacji zachowują orientację włókien.
Dla niezorientowanych
wiązek nad $S^1$ pozostają
dwie klasy w każdym dodatnim
rzędzie: wiązka trywialna
i wiązka z jednym skręceniem
typu Möbiusa.
\end{twierdzenie}
\begin{proof}
Wybierzmy metrykę na włóknach
i dodatnio zorientowane ramy
nad $D^m_+$ i $D^m_-$.
Na ich wspólnym równiku
ramy różnią się macierzą
$g(x)\in SO(r)$. Otrzymujemy
opisane sklejenie, więc
każda wiązka jest takiej
postaci. Wybieramy ramy gładkie, na przykład przez transport
wzdłuż promieni dysków. Wtedy $g$ jest gładka;
klasy izomorfizmu
topologicznego i gładkiego
na sferze są tu zgodne
przez gładkie przybliżenie.

Zmiana ramy na górnym
lub dolnym dysku mnoży
$g$ z lewej lub prawej
przez ograniczenie
$h_\pm:D^m_\pm\to SO(r)$.
Każde takie ograniczenie
jest homotopijne do stałej:
ściągamy dysk do punktu.
Ponieważ $SO(r)$ jest
spójna dla $r\geq2$,
stałe macierze można
dalej połączyć z
jednostkową. Zatem
izomorfizm wiązek
nie zmienia klasy
homotopii $g$.
Dla $r=1$ grupa $SO(1)$
jest trywialna i wniosek
również jest prawdziwy.

Odwrotnie, gdy $g_t$ jest homotopią funkcji
sklejających, wygładzamy ją względnie końców i sklejamy
wiązki trywialne nad
$D^m_\pm\times[0,1]$
za pomocą $g_t$.
Powstaje wiązka nad
$S^m\times[0,1]$.
Wybierzmy na niej
metrykę i koneksję
(lokalne zwykłe pochodne
można skleić podziałem
jedności). Transport
równoległy w kierunku
od $t=0$ do $t=1$
zależy gładko od punktu
sfery i daje izomorfizm
wiązek na obu końcach.
To dowodzi, że homotopijne
funkcje sklejające
dają izomorficzne wiązki.
Stałość punktu bazowego
nie zmienia klasyfikacji dla spójnej grupy $SO(r)$.
Nie tylko mapę, ale i homotopię $g_t$ można
znormalizować: $\widetilde g_t(x)=g_t(x_0)^{-1}g_t(x)$
stale przyjmuje wartość $I$ w $x_0$. To pokazuje
injektywność i surjektywność przejścia od swobodnych
do opartych klas homotopii.

Nad $S^1$ równik to
dwa punkty $S^0$.
Po zmianie ramy w
jednym z dwóch łuków
pozostaje tylko informacja,
czy transport dookoła
okręgu ma wyznacznik
$+1$ czy $-1$.
Składowa $GL_r(\R)$
o wyznaczniku dodatnim
jest spójna, podobnie
ujemna; otrzymujemy
dokładnie dwie klasy.
\end{proof}

\begin{figure}[htbp]
\centering
\begin{tikzpicture}[>=Latex,font=\small,line cap=round,scale=1.08]
 \shade[inner color=white,outer color=manifoldblue!43]
    (0,0) circle (1.55);
 \fill[white,opacity=.38] (-1.55,0)
     arc[start angle=180,end angle=360,x radius=1.55,y radius=.41]
     --cycle;
 \draw[manifoldblue!75,very thick]
     (-1.55,0)
     arc[start angle=180,end angle=360,x radius=1.55,y radius=.41];
 \draw[manifoldblue!55,dashed]
     (-1.55,0)
     arc[start angle=180,end angle=0,x radius=1.55,y radius=.41];
 \node[manifoldblue!80!black] at (0,1.05) {$D^m_+$};
 \node[manifoldblue!80!black] at (0,-1.05) {$D^m_-$};
 \draw[->,manifoldgold!85!black,very thick]
   (2.20,.54) .. controls (2.85,.55) and (2.85,-.55) .. (2.20,-.54);
 \node[anchor=west,manifoldgold!75!black] at (2.8,.03)
     {$g:S^{m-1}\to SO(r)$};
 \node[anchor=west] at (-1.30,-1.88)
     {dwie trywializacje spotykają się na równiku};
\end{tikzpicture}
\caption{Sfera jako dwa dyski z równikiem.
W każdej połowie wiązka jest iloczynem;
macierz $g(x)$ mówi, jak zidentyfikować
włókna po obu stronach równika.}
\label{fig:clutching-16}
\end{figure}

\begin{przyklad}[Dlaczego $TS^2$ jest skręcona]
Zorientowane wiązki płaszczyznowe
nad $S^2$ opisuje
$\pi_1(SO(2))=\pi_1(S^1)=\Z$.
Dla wiązki $TS^2$ wybierzmy
stereograficzne zespolone
współrzędne $z$ i $w$
na półkulach z przejściem
$w=1/z$ w pasie równika.
Pochodna przejścia
$\dd w/\dd z=-z^{-2}$
obraca ramę styczną
dwukrotnie, gdy $z$
okrąża równik raz.
Funkcja sklejająca ma
więc liczbę nawinięć
$\pm2$ (znak zależy
od wyboru orientacji
równika). Wiązka $TS^2$
nie jest trywialna:
na sferze nie istnieje
globalna rama dwóch
niezależnych pól stycznych.
\end{przyklad}

\begin{przyklad}[Grupa obrotów w wymiarze trzy]
Dla zorientowanych wiązek
rzędu $3$ nad $S^2$
otrzymujemy
$\pi_1(SO(3))\cong\Z/2$:
są dwie klasy.
Opis podwójnego nakrycia
$\operatorname{Spin}(3)\to SO(3)$
z Rozdziału~\ref{ch:grass-spin}
pokazuje, skąd bierze się
dwójka. Pętla macierzy
obrotu reprezentująca
klasę niezerową podnosi
się do drogi kończącej
w drugim punkcie włókna
nakrycia. Takiej
funkcji sklejającej
nie można podnieść
do zamkniętej pętli
w $\operatorname{Spin}(3)$;
wiązka ma niezerowe
$w_2\in H^2(S^2;\mathbb F_2)$.
Wiązka trywialna ma
$w_2=0$.
\end{przyklad}

\section{Klasy charakterystyczne: przeszkody w kohomologii}
\label{sec:klasy-char-16}

\begin{definicja}[Idea klasy charakterystycznej]
\index{idea klasy charakterystycznej@Idea klasy charakterystycznej}
Klasą charakterystyczną
wiązki $E\to M$ nazywamy
przypisanie $E\mapsto c(E)
\in H^k(M;G)$ zgodne
z cofnięciem:
$c(f^*E)=f^*c(E)$.
Jeżeli wiązka jest trywialna,
jej klasy dodatniego stopnia
są zerowe, bo jest cofnięciem
wiązki nad punktem.
Odwrotna implikacja
nie musi zachodzić:
klasy charakterystyczne
nie rozróżniają wszystkich
wiązek.
\end{definicja}

Najbardziej elementarna
przeszkoda dotyczy
orientacji. Nie trzeba
jej definiować przez
ogólną teorię przestrzeni
klasyfikujących.

\begin{stwierdzenie}[Pierwsza klasa Stiefela--Whitneya]
\label{prop:w1-orientacja-16}
Niech $E\to M$ będzie
rzeczywistą wiązką
rzędu $r$. Wybierzmy
trywializacje nad dobrym
pokryciem $\{U_i\}$ i
funkcje przejścia
$g_{ij}:U_i\cap U_j\to GL_r(\R)$.
Znaki
$\epsilon_{ij}=\operatorname{sgn}
\det g_{ij}\in\{\pm1\}$
wyznaczają klasę
$w_1(E)\in H^1(M;\mathbb F_2)$.
Mamy $w_1(E)=0$
wtedy i tylko wtedy,
gdy $E$ jest orientowalna.
Ponadto
$w_1(E\oplus F)=w_1(E)+w_1(F)$.
\end{stwierdzenie}
\begin{proof}
Na potrójnym przecięciu
$g_{ij}g_{jk}=g_{ik}$,
a więc
$\epsilon_{ij}\epsilon_{jk}
=\epsilon_{ik}$.
Po utożsamieniu $-1$
z $1\in\mathbb F_2$
jest to równanie jednokocyklu. Aby utożsamić jego
klasę z kohomologią singularną, można obyć się
bez dodatkowego formalizmu: mnożymy znaki przejść
wzdłuż pętli. Warunek na potrójnych przecięciach
zapewnia niezmienniczość przy podziale i homotopii
pętli, a więc homomorfizm $\pi_1(M)\to\mathbb F_2$.
Utożsamienia $H_1(M;\Z)=\pi_1(M)_{\mathrm{ab}}$
i $H^1(M;\mathbb F_2)=\operatorname{Hom}(H_1(M;\Z),\mathbb F_2)$
dają żądaną klasę na każdej składowej. Gdy homomorfizm
znika, transport orientacji z ustalonego punktu nie
zależy od drogi i nadaje zgodne orientacje lokalnym ramom.
Zmiana
ramy na $U_i$ mnoży
każdy znak $\epsilon_{ij}$
przez znak nowej ramy
na $U_i$ i na $U_j$;
odpowiada to dodaniu
kobrzegu, więc klasa
nie zależy od ram.
Jeśli $w_1(E)=0$,
kocykl znaków jest
kobrzegiem: zmieniamy
orientacje lokalnych
ram tak, aby wszystkie
macierze przejścia
miały dodatni wyznacznik.
Ramy nadają wtedy
spójną orientację
wiązce. Odwrotnie
spójna orientacja
pozwala od początku
wybrać takie ramy
i daje kocykel zerowy.
Macierz przejścia
dla $E\oplus F$
jest blokowa,
$\det(g_{ij}\oplus h_{ij})
=\det(g_{ij})\det(h_{ij})$;
znaki mnożą się, a
w $\mathbb F_2$
ich klasy się dodają.
\end{proof}

\begin{przyklad}
Wiązka Möbiusa nad $S^1$
ma monodromię $-1$,
więc $w_1$ jest
niezerowym elementem
$H^1(S^1;\mathbb F_2)$.
Trywialna linia ma $w_1=0$.
Dla rozmaitości
$w_1(TM)$ jest dokładnie
przeszkodą do jej
orientowalności.
Jest to również
monodromia systemu
$\mathcal O_M$
z Sekcji~\ref{sec:lokalne-wsp-16}.
\end{przyklad}

\begin{twierdzenie}[Pełna rodzina $w_i$]
\label{thm:wi-klasy-16}
Dla wiązki rzeczywistej
$E$ rzędu $r$ istnieją
klasy Stiefela--Whitneya
\[
 w_i(E)\in H^i(M;\mathbb F_2),
 \quad 0\leq i\leq r,\qquad
 w_0=1,\quad w_i=0\ (i>r),
\]
naturalne przy cofnięciu.
Dla sumy wiązek
\[
 w(E\oplus F)=w(E)\smile w(F),
 \qquad w(E)=1+w_1(E)+\cdots+w_r(E).
\]
W szczególności dla zanurzenia
$S\subset M$ z
Rozdziału~\ref{ch:wlozenia-otoczenia}
zachodzi
\begin{equation}\label{eq:w-normalne-16}
 w(TM|_S)=w(TS)\smile w(\nu(S,M)).
\end{equation}
\end{twierdzenie}
\begin{proof}
Dla $r=0$ przyjmujemy $w(E)=1$; dalej $r\geq1$.
Korzystamy z niezależnej konstrukcji klasy Thoma
i klasy $e_2$ poniżej. Dowód tej konstrukcji nie
zakłada istnienia $w_i$; zależności wyjaśnia uwaga
po Twierdzeniu~\ref{thm:euler-thom-16}.

Uzasadnijmy najpierw potrzebny rachunek we włóknie.
Na $\R P^m$ niech $x=w_1(\lambda)$ dla linii
tautologicznej. Ograniczenie ustalonego funkcjonału
liniowego do $\lambda$ jest przekrojem $\lambda^*$,
którego zera tworzą hiperpłaszczyznę $\R P^{m-1}$.
Parzystość przecięć pętli z tą hiperpłaszczyzną
jest zmianą znaku jej podniesienia do jednostkowych
wektorów linii: funkcjonał zmienia znak przy każdym
poprzecznym przejściu przez zero. Jest to dokładnie
monodromia $w_1$, więc hiperpłaszczyzna jest dualna do $x$.
$j$ niezależnych hiperpłaszczyzn przecina się
transwersalnie, a w $\R P^j$ ich przecięcie jest
jednym punktem. Modułowa wersja rachunku przecięć daje
$\langle x^j,[\R P^j]_2\rangle=1$.
Z obliczenia komórkowego grup $\R P^m$ wiemy, że
każda grupa $H^j(-;\mathbb F_2)$ dla $0\leq j\leq m$
jest jednowymiarowa. Zatem $1,x,\ldots,x^m$ są bazą,
a $x^{m+1}=0$ z przyczyn wymiarowych.

Podajemy konstrukcję algebraiczną klas ogólnej wiązki.
Niech $P(E)$ będzie wiązką
prostych we włóknach $E$,
$\lambda\to P(E)$ wiązką
tautologiczną, a
$x=w_1(\lambda)\in
H^1(P(E);\mathbb F_2)$.
Włóknem $P(E)$ jest
$\R P^{r-1}$; klasy
$1,x,\ldots,x^{r-1}$
po ograniczeniu do
każdego włókna tworzą
bazę jego kohomologii.
Stąd
\[
 H^*(P(E);\mathbb F_2)
  =\bigoplus_{j=0}^{r-1}
     H^*(M;\mathbb F_2)\,x^j
\]
jako moduł nad $H^*(M)$.
To twierdzenie o module
można sprawdzić komórka
po komórce: nad otwartą
komórką $E$ jest
trywialna, więc
$P(E)$ jest iloczynem
z $\R P^{r-1}$.
W względnym ciągu
dokładnym dołączania
komórki mapy
$1,x,\ldots,x^{r-1}$
są izomorfizmami na
względnych grupach
każdego stopnia.
Indukcja i lemat
o pięciu odwzorowaniach
dają powyższy rozkład
na całej skończonej
triangulacji. Dla
niezwartej bazy
ten sam rachunek
przeprowadza się
na filtracji przez
szkielety lokalnie
skończonego kompleksu:
relatywne grupy przy
dołączaniu wszystkich
komórek danego stopnia
są produktami grup
dla pojedynczych komórek,
a mapy pozostają
izomorfizmami.

Klasa $x^r$ ma wobec
tej bazy jednoznaczny
rozkład. Zapisujemy
go w postaci
\[
 x^r+
   \pi^*w_1(E)x^{r-1}
     +\cdots+\pi^*w_r(E)=0.
\]
Współczynnik $w_i$
ma stopień $i$,
bo każdy wyraz ma
całkowity stopień $r$.
Jednoznaczność rozkładu
i zgodność konstrukcji
$P(E)$ z cofnięciem
dają naturalność.
Przy rozszczepieniu
$E=L_1\oplus\cdots\oplus L_r$
relacja włóknowa
ma czynniki
$\prod_j(x+w_1(L_j))$.
Oto bezpośrednie
sprawdzenie. Nad $P(E)$
tautologiczna linia
$\lambda$ wchodzi
do $\pi^*E$, więc
homomorfizm inkluzji
jest nigdzie
nieznikającym
przekrojem wiązki
$\operatorname{Hom}(\lambda,\pi^*E)
=\bigoplus_j(\lambda^*\otimes\pi^*L_j)$.
Jej klasa Eulera modulo $2$, oznaczana niezależnie
przez $e_2$, jest zatem zerowa
na mocy
Twierdzenia~\ref{thm:euler-thom-16}
i uwagi o redukcji
modulo $2$ w
Sekcji~\ref{sec:klasy-eulera-16}.
Klasa Eulera modulo $2$
linii jest jej
$w_1$, a
$w_1(\lambda^*\otimes L_j)
=x+w_1(L_j)$
z rachunku znaków
funkcji przejścia.
Multiplikatywność
klasy Eulera sumy
daje właśnie
$\prod_j(x+w_1(L_j))=0$.
Porównanie współczynników
daje
$w(E)=\prod_j(1+w_1(L_j))$.

Każdą wiązkę można
rozszczepić po cofnięciu
do wiązki pełnych flag:
po kolei wybieramy
prostą w $E$ i w
ortogonalnym do niej
dopełnieniu. W każdym
takim kroku powyższy
rozkład modułowy
zawiera składnik $1$,
więc cofnięcie
na kohomologiach
jest injektywne.
Wzór iloczynowy
sprawdzony po
rozszczepieniu
schodzi zatem na $M$
i daje twierdzenie o sumie Whitneya. Dla dwóch
wiązek najpierw rozszczepiamy $E$, potem cofniętą $F$;
oba cofnięcia są injektywne, a obie wiązki rozszczepione.
Współczynnik stopnia jeden jest $\sum_jw_1(L_j)$,
czyli wcześniej zdefiniowany znak wyznacznika przejść.
Nowa definicja $w_1$ jest więc zgodna z poprzednią.
Wreszcie
$TM|_S\cong TS\oplus\nu$
z
Stwierdzenia~\ref{prop:normalna-rozszczepienie};
podstawienie do wzoru
sumy daje~\eqref{eq:w-normalne-16}.
\end{proof}

\begin{uwaga}[Gdzie już pojawiło się $w_2$?]
W Rozdziale~\ref{ch:grass-spin}
skonstruowaliśmy
$w_2(E)$ dla zorientowanej
wiązki jako przeszkodę
podniesienia funkcji
przejścia z $SO(r)$
do $\operatorname{Spin}(r)$.
Wyjaśnijmy zgodność z klasą stopnia dwa z
Twierdzenia~\ref{thm:wi-klasy-16}. Wystarczy ograniczyć
się do $2$-szkieletu bazy: restrykcja $H^2(M;\mathbb F_2)$
do tego szkieletu jest injektywna, ponieważ grupy
kołańcuchów i kobrzegi w stopniach $1,2$ są te same.
Na kompleksie wymiaru dwa wiązka rzędu co najmniej trzy
ma niezerowy przekrój: wybieramy go na wierzchołkach,
przedłużamy po krawędziach przez spójność sfery włókna,
a po $2$-komórkach przez jej jednospójność.
Powtarzając, rozszczepiamy zorientowaną wiązkę na
trywialną część i zorientowaną płaszczyznę $L$.
Wzór sumy daje $w_2(E)=w_2(L)=e(L)\bmod2$.
Wybierzmy zorientowaną ramę $L$ na $1$-szkielecie;
jest to możliwe dzięki spójności $SO(2)$.
Na $2$-komórce, względem tej ustalonej ramy brzegowej,
ostatnia liczba jest parzystością liczby nawinięć
funkcji $SO(2)$, bo to
stopień przeszkody przedłużenia niezerowego przekroju.
Ta sama parzystość jest przeszkodą podniesienia pętli
do $\operatorname{Spin}(2)\to SO(2)$, nakrycia $z\mapsto z^2$.
Po dodaniu kierunków trywialnych obrót o $2\pi$
nadal podnosi się do drogi od $1$ do $-1$ w grupie Spin.
Oba kocykle na $2$-komórkach są więc równe,
co dowodzi zgodności klas. Dla rzędu jeden obie są zerowe.
Z warunku
\eqref{eq:w-normalne-16}
wynika, że jeśli
$\nu(S,M)$ ma obramowanie,
to $w(\nu)=1$ i
$w(TM)|_S=w(TS)$.
To prosty test konieczny
dla niektórych zanurzeń
obramowanych.
\end{uwaga}

\begin{uwaga}[Klasy zespolone i Pontriagina]
Dla zespolonej wiązki
rzędu $r$ analogiczna
konstrukcja z
$\C P^{r-1}$
i klasą pierwszą
tautologicznej linii
o stopniu $2$
daje klasy Cherna
$c_i(E)\in H^{2i}(M;\Z)$.
Dla wiązki rzeczywistej
klasy Pontriagina
definiuje się przez
$p_i(E)=(-1)^i
c_{2i}(E\otimes_\R\C)
\in H^{4i}(M;\Z)$.
Ich naturalność i
zgodność ze stabilnym
dołożeniem wiązki
trywialnej wynikają
z konstrukcji projektowej.
W tym rozdziale
szczegółowo stosujemy
$w_1,w_2$ i klasę
Eulera; pełny rachunek
klas $c_i,p_i$
oraz konstrukcja
Cherna--Weila
wykraczają poza
niezbędne przygotowanie
do chirurgii.
\end{uwaga}

\section{Klasa Eulera, zera pól i samoprzecięcie}
\label{sec:klasy-eulera-16}

\begin{twierdzenie}[Klasa Thoma i klasa Eulera]
\label{thm:euler-thom-16}
Niech $E\to M$ będzie
zorientowaną rzeczywistą
wiązką rzędu $r$ z metryką
we włóknach. Dla jej
wiązki dysków $D(E)$
i wiązki sfer $S(E)$
istnieje klasa
\[
 u_E\in H^r(D(E),S(E);\Z),
\]
której ograniczenie do
każdego włókna
$(D^r,S^{r-1})$ jest
dodatnim generatorem.
Iloczyn z $u_E$ daje
izomorfizm Thoma
\[
 H^j(M;\Z)\stackrel{\cong}{\longrightarrow}
 H^{j+r}(D(E),S(E);\Z).
\]
Jeśli $s_0:M\to D(E)$
jest przekrojem zerowym,
to \emph{klasa Eulera}
$e(E)=s_0^*u_E\in H^r(M;\Z)$
jest naturalna przy
cofnięciu i spełnia
$e(E\oplus F)=e(E)\smile e(F)$
dla zorientowanych wiązek.
Istnienie nigdzie nieznikającego
przekroju pociąga $e(E)=0$.
\end{twierdzenie}

W zapisie $s_0^*u_E$ najpierw zapominamy o względności,
przechodząc z $H^r(D(E),S(E))$ do $H^r(D(E))$,
a potem cofamy klasę przez $s_0$. Równoważnym modelem
klasy Thoma jest $H^r(E,E\setminus0)$: wycięcie
i radialne cofnięcie włókien utożsamiają go z modelem
dyskowym. Ten model pozwala cofać klasę przez przekrój
jako klasę pary $(M,M\setminus Z(s))$.

\begin{proof}
Nad komórką bazy wiązka jest trywialna, zatem odpowiednia
para w przestrzeni całkowitej jest iloczynem tej komórki
z parą $(D^r,S^{r-1})$.
Włóknowa grupa względna
w stopniu $r$ jest $\Z$;
orientacja określa jej
dodatni generator.
Na przecięciu sąsiednich
trywializacji macierz
przejścia ma dodatni
wyznacznik, więc przenosi
ten generator w niego
samego. Sama zgodność lokalnych generatorów nie jest jeszcze
dowodem istnienia globalnej klasy. Skonstruujmy ją
w kompleksie komórkowym pary $(D(E),S(E))$.
Nad każdą $d$-komórką bazy dołączamy względną
komórkę $D^d\times D^r$ wymiaru $d+r$;
część $D^d\times S^{r-1}$ leży już w $S(E)$,
a $\partial D^d\times D^r$ nad wcześniejszym szkieletem.
Orientujemy iloczyn kolejnością „baza, włókno”.
Liczby incydencji są tymi z bazy pomnożonymi
przez stopień transportu orientacji włókna, czyli $+1$.
Kompleks względny kołańcuchów jest więc kompleksem
bazy przesuniętym o $r$. Stały kocykl $1$ w stopniu
zero bazy wyznacza kocykl stopnia $r$ pary,
którego klasa to $u_E$. Jest to jedyna klasa
o wymaganych ograniczeniach: pod tym przesunięciem
odpowiada funkcji na składowych bazy o wartości $1$.
Dla niezwartej bazy używamy produktów współczynników
po komórkach; indukcja po szkieletach ma skończenie
wiele etapów, gdyż rozmaitość ma skończony wymiar.

Teraz sprawdzamy, że mnożenie przez tę już
skonstruowaną klasę daje deklarowany izomorfizm:
na dołączanej $d$-komórce
relatywne grupy pary
wiązki są przesunięciem
o $r$ relatywnych grup
$(D^d,S^{d-1})$.
Mnożenie przez lokalny
generator jest tam
izomorfizmem. W długich
ciągach dokładnych
dla starego szkieletu
i nowej komórki wszystkie
kwadraty są naturalne;
indukcja i lemat o
pięciu odwzorowaniach
dają izomorfizm w
każdym stopniu.
Mnożenie ma dokładnie postać
$a\mapsto\pi^*a\smile u_E$; na każdej komórce
jest opisanym przesunięciem przez włóknowy generator.
Dla rzędu zero przyjmujemy $S(E)=\varnothing$
i $u_E=e(E)=1$.

Cofnięcie wiązki przenosi
orientacje i włóknowe
generatory, więc przenosi
$u_E$ na klasę Thoma
cofniętej wiązki. Po
ograniczeniu do
przekroju zerowego
daje to naturalność
$e(f^*E)=f^*e(E)$.
Dla sumy $E\oplus F$ możemy użyć dysku w normie
maksimum, czyli $D(E)\times_M D(F)$. Jego brzeg to
$(S(E)\times_M D(F))\cup(D(E)\times_M S(F))$.
Cofamy obie klasy Thoma przez rzuty i bierzemy
ich względny iloczyn kubkowy. W każdym włóknie
jest on dodatnim generatorem w orientacji „$E$, potem $F$”,
więc z jedyności jest klasą Thoma sumy.
Zmiana normy przez radialny homeomorfizm nie zmienia
klasy. Cofnięcie do zera daje
$e(E\oplus F)=e(E)\smile e(F)$.

Jeśli $s$ nigdzie
nie znika, po unormowaniu
jest odwzorowaniem
$M\to S(E)$.
Włączenie $S(E)\subset D(E)$
powoduje $s^*u_E=0$,
bo względna klasa
na $S(E)$ znika.
Przekroje $s_0$ i $s$
są homotopijne jako
odwzorowania $M\to D(E)$
przez $t\mapsto ts$.
Ich cofnięcia klasy
$u_E$ są więc równe;
$e(E)=0$.
\end{proof}

\begin{uwaga}[Wersja modulo $2$ i brak błędnego koła]
Powyższy dowód nie używa klas Stiefela--Whitneya.
Z $\mathbb F_2$ działa dla każdej wiązki, ponieważ
zmiana orientacji włókna nie zmienia generatora.
Oznaczmy otrzymaną klasę Eulera modulo $2$ przez $e_2(E)$.
Jest naturalna, multiplikatywna i znika przy istnieniu
niezerowego przekroju, z dokładnie tych samych dowodów.

Dla linii $L$ mamy $e_2(L)=w_1(L)$. Wybierzmy
niezerowe wartości przekroju w wierzchołkach triangulacji
i przedłużmy go ogólnie na krawędzie. Parzystość zer
na krawędzi wynosi $1$ dokładnie wtedy, gdy wartości
na końcach mają przeciwne znaki po transporcie
w trywializacji tej krawędzi. Jest to kocykl zmiany
orientacji, czyli $w_1(L)$; zmiana wartości przy
wierzchołku dodaje kobrzeg. Z drugiej strony te same
parzystości są lokalnymi stopniami cofnięcia klasy Thoma,
więc reprezentują $e_2(L)$.

Te własności $e_2$ wystarczają do dowodu wzoru
w Twierdzeniu~\ref{thm:wi-klasy-16}; nie zakłada on
z góry równości $e_2=w_r$. Dopiero po rozszczepieniu
$E$ na linie otrzymujemy
$e_2(E)=\prod_jw_1(L_j)=w_r(E)$.
Injektywność cofnięcia do wiązki flag przenosi
tę równość na bazę. Dla wiązki zorientowanej
redukcja całkowitej klasy Thoma jest klasą Thoma modulo $2$,
więc również $e(E)\bmod2=w_r(E)$.
\end{uwaga}

\begin{twierdzenie}[Zera przekroju reprezentują klasę Eulera]
\label{thm:euler-zeros-16}
Niech $M^n$ będzie
zamknięta i zorientowana,
a $E\to M$ zorientowana
rzędu $r\leq n$.
Dla przekroju $s$
transwersalnego do zera
zbiór $Z(s)$ jest
zorientowaną
$(n-r)$-podrozmaitością i
\[
 \operatorname{PD}^{-1}[Z(s)]=e(E).
\]
Jeśli $r=n$, zera są
skończone i
\begin{equation}\label{eq:euler-zeros-16}
 \langle e(E),[M]\rangle
   =\sum_{x\in Z(s)}
       \operatorname{sgn}\det(\D_xs).
\end{equation}
W bazach zgodnych z
orientacją znak nie
zależy od wyboru
trywializacji.
\end{twierdzenie}
\begin{proof}
Regularność wynika z Twierdzenia~\ref{thm:zera-przekroju}.
Pochodna pionowa w zerze daje izomorfizm
$\nu(Z(s),M)\cong E|_{Z(s)}$. Nadajemy normalnej
orientację przeniesioną z $E$, a $Z(s)$ orientujemy
warunkiem „normalna, potem styczna daje orientację $M$”.

W modelu $H^r(E,E\setminus0)$ cofnięcie przez $s$
jest klasą $s^*u_E\in H^r(M,M\setminus Z(s))$.
W małej trywializacji ograniczenie $s$ do normalnego
dysku ma w zerze odwracalną różniczkę o dodatnim
wyznaczniku, więc jego lokalny stopień jest $+1$.
Wycięcie i otoczenie tubularne identyfikują
$s^*u_E$ z klasą Thoma normalnej do $Z(s)$.
Jej obraz w kohomologii bezwzględnej jest klasą
dualną do $Z(s)$: na każdej normalnej komórce
ocenia $+1$, zgodnie z orientacjami w dowodzie dualności.
Liniowa homotopia $s\simeq s_0$ w $E$ pokazuje,
że ten sam obraz jest $e(E)$.

Dla $r=n$ każdy punkt zerowy jest zorientowaną
$0$-komórką o współczynniku
$\operatorname{sgn}\det\D_xs$. Zwartość i izolacja
zer dają ich skończoność. Ewaluacja na klasie
fundamentalnej sumuje te współczynniki,
co daje~\eqref{eq:euler-zeros-16}.
\end{proof}

\begin{wniosek}[Samoprzecięcie]
\label{cor:samoprzeciecie-16}
Niech $M^{2k}$ będzie zorientowana, a
$S^k\subset M$ zamknięta i zorientowana.
W tym wzorze nadajemy wiązce normalnej orientację
przez kolejność \emph{styczna do $S$, potem normalna},
zgodną z orientacją $M$. Wówczas
to liczba przecięcia
$S$ z małą przesuniętą
kopią wynosi
\[
 I(S,S)=\langle e(\nu(S,M)),[S]\rangle.
\]
\end{wniosek}
\begin{proof}
W otoczeniu tubularnym
zapisujemy $S$ jako
przekrój zerowy $\nu$.
Wybierzmy mały przekrój
$s$ transwersalny do
zera i przesuńmy $S$
do jego wykresu.
Wykres przecina
przekrój zerowy
dokładnie dla $s(x)=0$.
W bazach „styczna, normalna” kolumny baz
przekroju zerowego i wykresu tworzą macierz
$\begin{psmallmatrix}I&I\\0&\D_xs\end{psmallmatrix}$.
Jej wyznacznik to $\det\D_xs$, więc znak przecięcia
jest lokalnym stopniem $s$. Ten rachunek ustala
potrzebną orientację normalnej; nie można jej
odwrócić dowolnie, zachowując ten sam zapis wzoru.
Suma znaków to
$\langle e(\nu),[S]\rangle$
z~\eqref{eq:euler-zeros-16}.
\end{proof}

\begin{przyklad}
Dla $S^2\times\{p\}
\subset S^2\times S^2$
wiązka normalna jest
trywialna: jej włókno
to stała przestrzeń
$T_pS^2$. Istnieje
niezerowy stały
przekrój, więc
$I(S,S)=0$.
Dwie różne sfery
$S^2\times\{p\}$
i $\{q\}\times S^2$
przecinają się
jednak w jednym
punkcie z liczbą
przecięcia $1$.
\end{przyklad}

\section{Twierdzenie Poincarégo--Hopfa}
\label{sec:poincare-hopf-16}

\begin{definicja}[Indeks zera pola]
\index{indeks zera pola@Indeks zera pola}
Jeżeli pole $X$ na
$n$-rozmaitości ma
izolowane zero $x$,
wybierzmy małą kulę w mapie wokół $x$
oraz trywializację stycznej przez różniczkę tej mapy.
\emph{Indeks} $\operatorname{ind}_x X$
jest stopniem mapy
$X/|X|:S^{n-1}\to S^{n-1}$, dla $n\geq1$.
Zmiana mapy zmienia jednocześnie orientację dziedziny
i współrzędnych wektora, więc dwa znaki się mnożą
do $+1$. Zmienną macierz przejścia na małej kuli
można zdeformować do jej wartości w środku, nie
przechodząc przez macierze osobliwe. Indeks nie
zależy więc od mapy ani od lokalnej orientacji.
Dla $n=1$ stopień liczymy na $\widetilde H_0(S^0)$;
dla $n=0$ umownie każdy punkt zerowy ma indeks $1$.
Jeśli $\D_xX$ jest
odwracalna,
$\operatorname{ind}_xX
=\operatorname{sgn}\det\D_xX$:
na małej sferze
$X(y)=\D_xX(y-x)+o(|y-x|)$
jest homotopijne
do odwracalnej
mapy liniowej.
\end{definicja}

\begin{twierdzenie}[Poincaré--Hopf]
\label{thm:poincare-hopf-16}
Jeśli $M$ jest
zamkniętą
zorientowaną rozmaitością,
a pole $X$ ma
izolowane zera, to
\begin{equation}\label{eq:poincare-hopf-16}
 \sum_{x:X(x)=0}
     \operatorname{ind}_x X
 =\langle e(TM),[M]\rangle
 =\chi(M).
\end{equation}
W szczególności
$\chi(M)\ne0$
uniemożliwia istnienie
nigdzie nieznikającego
pola stycznego.
\end{twierdzenie}
\begin{proof}
Dla $n=0$ obie strony liczą punkty $M$.
Dalej przyjmujemy $n\geq1$.
\emph{Od indeksów do klasy Eulera.}
Przy zerze nieosobliwym
pierwsza równość jest
dokładnie
Twierdzeniem~\ref{thm:euler-zeros-16}
dla $E=TM$. Dla zera
izolowanego, lecz
osobliwego, wybieramy
kulę bez zer na
brzegu. Mała perturbacja
pola wewnątrz kuli,
niezmieniająca go
przy brzegu, może być
transwersalna do
przekroju zerowego
na mocy
Rozdziału~\ref{ch:sard-transwersalnosc}.
Stopień mapy
$X/|X|$ na brzegu
nie zmienia się
w tej homotopii;
równa się sumie
znaków nowo
powstałych zer
wewnątrz. Zwartość
$M$ pozwala wykonać
ten zabieg w
skończenie wielu
rozłącznych kulach.

\emph{Dlaczego suma jest $\chi(M)$?}
Przypomnijmy i uzasadnijmy
potrzebną postać
zasady Lefschetza.
Dla gładkiego
$f:M\to M$ o
skończenie wielu
nieosobliwych
punktach stałych
zachodzi
\begin{equation}\label{eq:lefschetz-lokalny-16}
 \sum_{f(x)=x}
    \operatorname{sgn}\det(I-\D_xf)
 =\sum_j(-1)^j
    \operatorname{tr}
      (f_*:H_j(M;\Q)\to H_j(M;\Q)).
\end{equation}
Nadajemy grafowi $\Gamma_f$ i przekątnej $\Delta$
orientacje przeniesione z $M$, a $M\times M$ orientację
produktową. W punkcie stałym, w tej kolejności płatów,
macierz baz stycznych wynosi
\[
 \begin{pmatrix}I&I\\\D_xf&I\end{pmatrix},
 \qquad \det=\det(I-\D_xf).
\]
Lewa strona~\eqref{eq:lefschetz-lokalny-16}
jest więc $I(\Gamma_f,\Delta)$.

Obliczymy tę liczbę bez pomijania znaku globalnego.
Wybierzmy bazę $a_{ij}$ przestrzeni $H^i(M;\Q)$
i bazę dualną $b_{ij}\in H^{n-i}(M;\Q)$, tzn.
$\langle a_{ij}\smile b_{i\ell},[M]\rangle=\delta_{j\ell}$.
Przez $u\times v$ oznaczamy
$\pr_1^*u\smile\pr_2^*v$. Te iloczyny tworzą bazę
kohomologii produktu nad $\Q$. W tym szczególnym
przypadku formułę Künnetha można uzasadnić bez torsji:
komórki produktu mają brzeg
$\partial(e^p\times e^q)=\partial e^p\times e^q+
(-1)^p e^p\times\partial e^q$.
Ich kompleks jest tensorem kompleksów czynników.
Nad ciałem każdy skończony kompleks rozkłada się na
homologię z zerową różniczką i dwuwyrazowe kompleksy
izomorfizmów; tensor z każdym z tych ostatnich
jest ściągalny. Pozostaje tensor homologii,
a po dualizacji podana baza kohomologii.

Klasa dualna przekątnej ma postać
\[
 \alpha_\Delta=\sum_{i,j}(-1)^{ni}\,b_{ij}\times a_{ij}.
\]
Sprawdzamy ją przez ewaluację na bazowych iloczynach:
z definicji dualności
$\langle\alpha_\Delta\smile(u\times v),[M\times M]\rangle
=\langle u\smile v,[M]\rangle$.
Dla $|u|=i$, $|v|=n-i$ znak przestawienia czynników
w iloczynie $(b_{ij}\times a_{ij})\smile(u\times v)$
wynosi $(-1)^{i^2}$, a zamiana kolejności $b_{ij},u$
daje $(-1)^{i(n-i)}$. Razem z $(-1)^{ni}$ jest to
$+1$. Dualność baz daje żądaną ewaluację; w innych
stopniach wkład jest zerowy. Niezdegenerowanie
parowania w produkcie dowodzi wzoru na $\alpha_\Delta$.

Teraz ograniczamy tę klasę do grafu, gdzie
$\pr_1$ jest identycznością, a $\pr_2=f$:
\begin{align*}
 I(\Gamma_f,\Delta)
 &=\langle\alpha_\Delta,[\Gamma_f]\rangle\\
 &=\sum_{i,j}(-1)^{ni}
       \langle b_{ij}\smile f^*a_{ij},[M]\rangle\\
 &=\sum_i(-1)^i\operatorname{tr}
       (f^*:H^i(M;\Q)\to H^i(M;\Q)).
\end{align*}
Ostatni znak wynika z
$ni+i(n-i)\equiv i\pmod2$. Sumowanie po $j$
wybiera elementy diagonalne macierzy $f^*$.
Ewaluacja identyfikuje $f^*$ z dualnym odwzorowaniem
do $f_*$, więc ich ślady są równe. Otrzymujemy
\eqref{eq:lefschetz-lokalny-16}.

Teraz przyjmijmy, że
zera $X$ są nieosobliwe,
i niech $f=\Phi_{-t}$
będzie przepływem
$X$ przez bardzo
mały czas $-t<0$.
Punkty zerowe są
punktami stałymi.
Poza małymi kulami
otaczającymi zera
norma $X$ ma
dodatnie minimum,
więc z rozwinięcia
$\Phi_{-t}(y)=y-tX(y)+O(t^2)$
jednolitego na tym
zbiorze nie ma
innych punktów
stałych dla małego $t$.
W mapie przy zerze wyrażenie
$F(y,t)=(y-\Phi_{-t}(y))/t$ przedłuża się gładko
do $t=0$ przez $F(y,0)=X(y)$; wynika to ze wzoru
całkowego na przyrost przepływu. Ponieważ
$\D_yF(x,0)=\D_xX$ jest odwracalna, twierdzenie
o funkcji uwikłanej z parametrem $t$ daje jedyne
rozwiązanie w jednym ustalonym otoczeniu $x$.
Jest nim stale $x$. Zmniejszamy kule do tych
otoczeń, a ich zwarte dopełnienie obejmuje także
pozostałe pierścienie wokół zer. Poprzednie
jednolite oszacowanie wyklucza tam punkty stałe.
W punkcie zerowym
\[
 I-\D_x\Phi_{-t}
   =t\D_xX+O(t^2),
\]
zatem znak wyznacznika
to $\operatorname{ind}_xX$.
Przepływ jest
homotopijny do
identyczności,
więc $f_*$ jest
identycznością
na wszystkich
$H_j(M;\Q)$.
Prawa strona
\eqref{eq:lefschetz-lokalny-16}
wynosi
$\sum_j(-1)^j b_j=\chi(M)$.
Otrzymujemy drugą
równość~\eqref{eq:poincare-hopf-16}.
Perturbacja z początku
dowodu rozszerza ją
na zera osobliwe.
\end{proof}

\begin{figure}[htbp]
\centering
\begin{tikzpicture}[>=Latex,font=\small,line cap=round,scale=1.05]
 \shade[inner color=white,outer color=manifoldblue!43] (0,0) circle (1.55);
 \foreach \c in {-.88,-.48,0,.48,.88}
  \draw[manifoldblue!55,thin]
   plot[domain=8:172,samples=65,variable=\t]
    ({1.55*\c*sin(\t)},{1.55*cos(\t)});
 \foreach \c in {-.88,-.48,0,.48,.88}
  \draw[->,manifoldgold!85!black,thick]
   plot[domain=115:62,samples=25,variable=\t]
    ({1.55*\c*sin(\t)},{1.55*cos(\t)});
 \fill[manifoldred] (0,1.55) circle (2.7pt);
 \fill[manifoldred] (0,-1.55) circle (2.7pt);
 \node[anchor=west,manifoldred] at (.16,1.55) {$+1$};
 \node[anchor=west,manifoldred] at (.16,-1.55) {$+1$};
 \node[anchor=west] at (1.9,.30) {$h(x,y,z)=z$};
 \node[anchor=west] at (1.9,-.12) {$X=e_z-z(x,y,z)$};
\end{tikzpicture}
\caption{Gradient wysokości na $S^2$
ma dwa zera: na biegunie dolnym
i górnym. Oba mają indeks $+1$,
więc suma indeksów i $\chi(S^2)$ są równe $2$.
Złote łuki pokazują rzuty trajektorii gradientu
w metryce standardowej: biegną po południkach od minimum
do maksimum. Promień rysunku jest tylko skalą graficzną.}
\label{fig:poincare-hopf-16}
\end{figure}

\begin{przyklad}[Sfera i torus]
Na $S^2$ gradient
funkcji wysokości
ma minimum i maksimum.
Lokalne macierze
pochodnej pola są
odpowiednio dodatnio
i ujemnie określone
w dwóch wymiarach,
ale oba wyznaczniki
są dodatnie:
$1+1=\chi(S^2)=2$.
Na torusie $T^2$
istnieje nigdzie
nieznikające pole
styczne wzdłuż
jednego czynnika
$S^1$, więc
$e(TT^2)=0$ i
$\chi(T^2)=0$.
To dopowiada topologiczną
treść wzoru
Gaussa--Bonneta z
Rozdziału~\ref{ch:krzywizna}.
\end{przyklad}

\section{Forma przecięcia i przygotowanie do chirurgii}
\label{sec:forma-przeciecia-16}

Gdy $n=2k$, dwie
$k$-podrozmaitości
przecinają się
w punktach.
Ich liczba przecięcia
tworzy na środkowej
homologii formę
dwuliniową. Część
torsyjna nie jest
wykrywana przez
wartości całkowite
tej formy, więc
najpierw ją usuwamy.

\begin{twierdzenie}[Forma przecięcia]
\label{thm:forma-przeciecia-16}
Dla zamkniętej
zorientowanej
$M^{2k}$ parowanie
\[
 Q_M:
 \frac{H_k(M;\Z)}{\mathrm{Tor}}
 \times
 \frac{H_k(M;\Z)}{\mathrm{Tor}}
   \longrightarrow\Z,
 \qquad
 Q_M(a,b)=
 \langle\operatorname{PD}^{-1}(a)
 \smile\operatorname{PD}^{-1}(b),[M]\rangle
\]
jest \emph{unimodularne}, czyli doskonałe nad $\Z$:
macierz w dowolnej bazie
części wolnej ma wyznacznik
$\pm1$. Ponadto
$Q_M(a,b)=(-1)^kQ_M(b,a)$.
W wymiarze $4m$ forma
jest symetryczna, a
w wymiarze $4m+2$
antysymetryczna.
\end{twierdzenie}
\begin{proof}
Z dualności Poincarégo
otrzymujemy
$H_k(M;\Z)\cong H^k(M;\Z)$.
Twierdzenie o
współczynnikach
uniwersalnych daje
ciąg dokładny
\[
0\longrightarrow
 \operatorname{Ext}
     (H_{k-1}(M),\Z)
 \longrightarrow H^k(M;\Z)
 \longrightarrow
 \operatorname{Hom}(H_k(M),\Z)
 \longrightarrow0.
\]
Grupy homologii
zwartej triangulacji
są skończenie
generowane.
Składnik $\operatorname{Ext}$
jest więc torsyjny,
a grupa $\operatorname{Hom}$
jest dualną kratą
części wolnej $H_k$.
Po podzieleniu przez
torsję otrzymujemy
izomorfizm
\[
 H_k(M)/\mathrm{Tor}
 \ \cong\
 \operatorname{Hom}
    (H_k(M)/\mathrm{Tor},\Z).
\]
Dla wybranej konwencji kapowej ten izomorfizm
wysyła $a$ na funkcjonał
\[
 b\longmapsto\langle\operatorname{PD}^{-1}(a),b\rangle
 =Q_M(b,a)=(-1)^kQ_M(a,b).
\]
Wspólny czynnik $(-1)^k$ nie zmienia odwracalności.
Klasy torsyjne nie wpływają na $Q_M$, ponieważ
ich całkowita ewaluacja jest torsyjną liczbą całkowitą,
a więc zerem. Definicja na ilorazie jest zatem jednoznaczna.
Macierz odwzorowania
między wolną grupą
abelową a jej dualną
jest odwracalna nad
$\Z$, stąd jej
wyznacznik to $\pm1$.
Dla wszystkich klas, także niereprezentowanych
przez zanurzone podrozmaitości, przemienność
stopniowana iloczynu kubkowego daje znak
$(-1)^{k^2}=(-1)^k$. To dowodzi obu symetrii.
Nie zakładamy, że każda całkowita klasa homologii
ma reprezentanta będącego podrozmaitością.
\end{proof}

\begin{przyklad}
Na $S^2\times S^2$
w bazie
$[S^2\times\{p\}]$,
$[\{q\}\times S^2]$
forma ma macierz
\[
 Q=
 \begin{pmatrix}
  0&1\\1&0
 \end{pmatrix},
 \qquad\det Q=-1.
\]
Na torusie
$H_1(T^2)=\Z^2$
i w bazie dwóch
podstawowych okręgów
otrzymujemy
$\begin{psmallmatrix}
0&1\\-1&0
\end{psmallmatrix}$.
Druga forma jest
antysymetryczna,
bo oba cykle mają
wymiar nieparzysty.
\end{przyklad}

\begin{definicja}[Elementarna operacja chirurgiczna]
\index{elementarna operacja chirurgiczna@Elementarna operacja chirurgiczna}
\label{def:chirurgia-zarys-16}
Niech $0\leq k\leq n-1$
i niech dane będzie
\emph{obramowane}
zanurzenie
$S^k\times D^{\,n-k}
\hookrightarrow M^n$.
Elementarna
\emph{chirurgia indeksu $k$} (w przyjętej tutaj
konwencji: wymiar operowanej sfery; uchwyt w jej
śladzie ma indeks $k+1$) polega na usunięciu
wnętrza tego
otoczenia i
wklejeniu
$D^{k+1}\times
S^{\,n-k-1}$
wzdłuż wspólnego brzegu
\[
 \partial(S^k\times D^{n-k})
 =S^k\times S^{n-k-1}
 =\partial(D^{k+1}\times S^{n-k-1}).
\]
Dla ustalonego obramowanego zanurzenia sklejamy
punkty o tych samych współrzędnych produktowych na
$S^k\times S^{n-k-1}$; jest to określony dyfeomorfizm
brzegów. Jego dowolna dodatkowa zmiana jest dodatkową
daną i może zmienić wynik. Kołnierze określają gładką
strukturę przy szwie. Jeśli chcemy zachować orientację,
identyfikacja musi odwracać orientacje indukowane
na brzegach obu sklejanych części.
\end{definicja}

To wyjaśnia rolę
wcześniejszych tematów.
Twierdzenie o otoczeniu
tubularnym zapewnia
obszar, który można
wyciąć. Obramowanie
z Definicji~\ref{def:obramowanie-normalne}
pozwala zapisać go
jako \emph{iloczyn}
$S^k\times D^{n-k}$.
Klasy charakterystyczne
mogą wykazać, że
normalna wiązka
nie ma takiego
obramowania.
Forma przecięcia
kontroluje, jak
cykle środkowego
wymiaru są ze sobą
powiązane; po
operacji trzeba
obliczyć ją na nowo.
Szczegółowa teoria
sklejania, uchwytów
i zmian homologii
będzie osobnym
rozdziałem.

\begin{przyklad}[Cięcie sfery w dwóch punktach]
Na $S^2$ wybierzmy
dwa rozłączne
małe dyski:
$S^0\times D^2$
to właśnie ich suma.
Po wycięciu wnętrz
pozostaje walec
$S^1\times[0,1]$.
W chirurgii
$k=0,n=2$
dołączamy
$D^1\times S^1$,
drugi walec,
wzdłuż dwóch
brzegowych okręgów.
Wybieramy sklejenia zgodne z orientacją powstałej
powierzchni, czyli odwracające orientacje jej brzegów.
Wtedy dwa walce składają się w zorientowaną zamkniętą
powierzchnię rodzaju jeden, czyli torus. Bez tego
warunku sklejenie może dać butelkę Kleina:
po obejściu podłużnego okręgu kierunek poprzeczny
wraca wówczas z przeciwną orientacją. Sama
charakterystyka Eulera nie rozróżnia tych wyników.
Na poziomie
charakterystyki
Eulera: ze sfery
o $\chi=2$ usunięcie
dwóch dysków
odejmuje $2$,
a dołączenie walca nie zmienia wyniku:
$\chi(A\cup B)=\chi(A)+\chi(B)-\chi(A\cap B)$,
przy czym zarówno walec, jak i dwa okręgi sklejenia
mają charakterystykę $0$.
Torus ma więc
$\chi=0$, zgodnie
z bezpośrednim
obliczeniem.
\end{przyklad}

\par\smallskip
{\footnotesize\noindent\emph{Dalsza lektura:}
\href{https://pi.math.cornell.edu/~hatcher/AT/AT.pdf}{A.\ Hatcher, \emph{Algebraic Topology}, dualność i iloczyn kubkowy},
\href{https://pi.math.cornell.edu/~hatcher/VBKT/VB.pdf}{A.\ Hatcher, \emph{Vector Bundles and K-Theory}, sklejanie i klasy charakterystyczne},
\href{https://math.stanford.edu/~ralph/bookR4.pdf}{R.\ Cohen, \emph{Bundles, Manifolds, and Homotopy}, przecięcia i klasa Thoma}.\par}
